\section{An integral inequality for convex tubes}\label{sec:tubes}

Let $E\subset\R^{n+1}$ be a closed full-dimensional convex set whose
boundary is smooth, has positive second fundamental form, and is
stationary for affine area.  In the applications, $E$ will be an affine
image of the original epigraph.  Fix a splitting
\[
 \R^{n+1}=\R^k_s\times\R^m_y,
 \qquad 2\le k\le n,\qquad m=n+1-k,
\]
and an open box $\mathcal B\subset(0,\infty)^k$.  Suppose that for every
$s\in\mathcal B$ the fiber
\[
 K_s=\{y\in\R^m:(s,y)\in E\}
\]
is compact and contains $0$ in its interior.  Our goal is the logarithmic
inequality in Proposition~\ref{prop:log-inequality}, whose constant depends
only on $n$, independently of the box and the shapes of its fibers.

\subsection{Support coordinates and invariant measures}

Fibers over compact subsets of $\mathcal B$ are uniformly bounded.
Otherwise there would be $(s_j,y_j)\in E$ with $s_j$ bounded and
$|y_j|\to\infty$.  After passage to a subsequence,
$y_j/|y_j|\to e\in S^{m-1}$.  Fix $s_0\in\mathcal B$, so that
$(s_0,0)\in E$.  For every fixed $t\ge0$, convexity and closedness give
\[
 (s_0,te)=\lim_{j\to\infty}
 \left[\left(1-\frac{t}{|y_j|}\right)(s_0,0)
       +\frac{t}{|y_j|}(s_j,y_j)\right]\in E.
\]
This contradicts compactness of $K_{s_0}$.  Thus the part of $E$ over
any compact base subset is compact.

The projection of $\partial E$ to the $s$ variables has full rank above
$\mathcal B$.  Indeed, otherwise a conormal at such a point would have
zero $y$ component.  Its supporting hyperplane would then give a
supporting hyperplane to the projection of $E$ at an interior point of
that projection, which is impossible.  Thus $\partial K_s$ is smooth,
and the restriction of the second fundamental form makes it strictly
positively curved.  Its outward Gauss map is a diffeomorphism onto
$S^{m-1}$ when $m>1$.  The Gauss maps also depend smoothly on $s$, since
their angular differentials are invertible.  Consequently the support
function $h(s,e)=\max_{y\in K_s}e\cdot y$ is smooth and positive on
\[
 \mathcal T=\mathcal B\times S^{m-1},
\]
and the corresponding part of $\partial E$ is parametrized by
\begin{equation}\label{eq:tube-parametrization}
 X(s,e)=(s,Y(s,e)),
 \qquad Y=he+\nabla_S h.
\end{equation}
Here $\nabla_S$ is the round connection. This support-function
reconstruction and its angular matrix are classical; see, for example,
\cite[preprint version, Section~7, Equations~(7.15)--(7.17)]{TW2005}.
We compute below
how they enter the second form and affine-area density of the present
tube. When $m=1$, we regard
$S^0=\{-1,1\}$ as the two outward directions of the interval $K_s$;
Equation~\eqref{eq:tube-parametrization} then means $Y=he$ on the two
endpoint sheets.  Throughout this section, angular matrices in this
case have size zero, their determinants equal $1$, and angular
derivatives, traces, and contractions equal $0$.  The measure
$\dd\omega$ on $S^0$ is counting measure.

Set
\[
 B=-D_s^2h,
 \qquad R=\nabla_S^2h+h\Id,
 \qquad \nu=(D_sh,-e).
\]
The identities $e\cdot Y=h$, $e\cdot Y_s=h_s$, and
$e\cdot\dd_eY=0$ show that $\nu$ annihilates the tangent vectors of
$X$.  It is inward because the fiber lies on the side
$e\cdot y\le h(s,e)$.  Differentiating the same identities gives
\[
 \nu X_{s_i s_j}=-h_{s_i s_j},
 \qquad \nu X_{s_i e_a}=0,
 \qquad \nu X_{e_a e_b}=R_{ab}.
\]
For the last equality we used $\dd_eY=R$ in round orthonormal frames.
Thus the second form for this conormal is
\begin{equation}\label{eq:tube-second-form}
 H=\operatorname{diag}(B,R).
\end{equation}
Both blocks are positive definite.  Taking $\xi=(0,-e)$ gives
$\nu\xi=1$ and
$|\det(X_1,\ldots,X_n,\xi)|=\det R$.  The affine-area formula therefore
becomes
\begin{equation}\label{eq:tube-area}
 \dd A=I_0\,\dd s\,\dd\omega,
 \qquad I_0=(\det B)^\delta(\det R)^{1-\delta},
\end{equation}
where we used $(n+1)\delta=1-\delta$.

To compare fibers after linear normalizations, we also need a measure
invariant under separate linear changes of $s$ and $y$.
The vector $V=(0,-Y)$ points from the fiber boundary toward its origin,
and $\nu V=h>0$.  Normalizing the conormal to take value $1$ on $V$
therefore defines the positive metric
\[
 G=H/h.
\]
Its Riemannian volume is
\begin{equation}\label{eq:tube-volume}
 \dd\mu=\mu_0\,\dd s\,\dd\omega,
 \qquad
 \mu_0=h^{-n/2}\sqrt{\det B\,\det R}.
\end{equation}
This definition makes $G$ and $\mu$ invariant under separate invertible
linear changes of $s$ and $y$.  To check the normalization explicitly,
let $(s',y')=(As,Ly)$ and, for an old unit normal $e$, put
\[
 t=|L^{-\mathsf T}e|,
 \qquad e'=L^{-\mathsf T}e/t.
\]
The new support function obeys $h'(As,e')=h(s,e)/t$, and its conormal
is $\operatorname{diag}(A,L)^{-\mathsf T}\nu/t$.  Under the induced
parameter map the new second form pulls back to $H/t$.  Dividing it by
$h'=h/t$ recovers $G$, including the change in angular parameters.

For $s$ covectors we henceforth use the inner product and norm given by
$hB^{-1}$, and for angular covectors those given by $hR^{-1}$.  We write
both as $\langle\cdot,\cdot\rangle$ and $|\cdot|$, with the variables
specifying the block.  Thus, for example,
\[
 |D_sg|^2=h\,(D_sg)^{\mathsf T}B^{-1}D_sg.
\]
All unqualified integrals below are over $\mathcal T$.  Every cutoff in
the $s$ variables has compact support in $\mathcal B$, so these
integrals and the corresponding variations have compact support on
the hypersurface.

\subsection{The stationarity equation in support coordinates}

The support coordinates have identified both affine area and the
invariant volume. We next express stationarity through their density
ratio. This supplies the differential identity whose weighted integral
will have a strictly positive quadratic part in dimensions three
through nine.

We first record the cofactor identities needed for integration by
parts.  Since $B=D_s^2(-h)$, its cofactor is divergence free in the
flat variables.  The angular tensor $R$ is Codazzi:
\[
 \nabla_aR_{bc}=\nabla_bR_{ac}.
\]
One way to see this is to differentiate $\dd_eY=R$ in the ambient
Euclidean space.  Commuting the two derivatives of $Y$ and taking
tangential components gives the displayed identity; the normal
components cancel because $R$ is symmetric.  Equivalently, the
curvature terms from commuting third derivatives of $h$ cancel the
derivatives of $h\Id$.

For completeness, if $A$ is any symmetric Codazzi tensor of positive
size $r$, its cofactor is divergence free.  In normal orthonormal
frames the cofactor is
\[
 (\cof A)^{ij}
 =\frac{1}{(r-1)!}
   \delta^{i i_2\cdots i_r}_{j j_2\cdots j_r}
   A_{i_2j_2}\cdots A_{i_rj_r},
\]
where the symbol on the right is the alternating Kronecker symbol.
On differentiating and contracting the derivative index with $i$,
each resulting term contains
$\nabla_iA_{i_a j_a}=\nabla_{i_a}A_{i j_a}$, symmetric in two indices
against which that symbol is alternating.  Each term is therefore
zero.  Symmetry gives the divergence identity in either contracted
index.  For $r=1$ the cofactor is the constant $1$; for the empty
angular block no identity is needed.  In particular,
\begin{equation}\label{eq:tube-cofactor-divergence}
 \partial_{s_i}(\cof B)^{ij}=0,
 \qquad \nabla_a(\cof R)^{ab}=0.
\end{equation}

Define
\begin{equation}\label{eq:tube-af-definitions}
 a=\left(\frac{\det B}{\det R}\right)^\delta,
 \qquad f=\log(a/h),
 \qquad c=\frac{n+2}{2}.
\end{equation}
The function $f$ records the ratio of affine area to invariant volume:
by Equations~\eqref{eq:tube-area} and~\eqref{eq:tube-volume},
\[
 \frac{\dd A}{\dd\mu}
 =h^{n/2}\left(\frac{\det B}{\det R}\right)^{-n/(2(n+2))}
 =\left(\frac ah\right)^{-n/2}=e^{-nf/2}.
\]
A variation $h+\varepsilon\eta$, with
$\eta\in C_c^\infty(\mathcal T)$, gives a compactly supported smooth
variation of $X$.  Positivity persists for sufficiently small
$\varepsilon$ on its compact support.  Differentiating
Equation~\eqref{eq:tube-area}, using stationarity, and dividing by
$\delta$ gives
\[
 0=\int\left[
 -I_0 B^{-1}:D_s^2\eta
 +(n+1)I_0 R^{-1}:(\nabla_S^2\eta+\eta\Id)
 \right]\dd s\,\dd\omega.
\]
The colon denotes contraction in the indicated orthonormal frames.
The scalar factors in this expression satisfy
\[
 I_0B^{-1}=(\cof B)a^{-(n+1)},
 \qquad I_0R^{-1}=(\cof R)a.
\]
Two integrations by parts using
Equation~\eqref{eq:tube-cofactor-divergence} therefore yield
\begin{equation}\label{eq:tube-euler}
 -(\cof B):D_s^2(a^{-(n+1)})
 +(n+1)(\cof R):(\nabla_S^2a+a\Id)=0.
\end{equation}
There is no angular boundary term.  In the case $m=1$, the calculation
is performed separately on the two endpoint sheets and the second
term is absent.

Writing $t=\log a$, we have
\[
 \begin{split}
 D_s^2(a^{-(n+1)})
   &=(n+1)a^{-(n+1)}
     \bigl(-D_s^2t+(n+1)D_st\otimes D_st\bigr),\\
 \nabla_S^2a&=a\bigl(\nabla_S^2t+
                        \nabla_St\otimes\nabla_St\bigr).
 \end{split}
\]
Dividing Equation~\eqref{eq:tube-euler} by $(n+1)I_0$ gives
\begin{equation}\label{eq:tube-log-euler}
 \begin{split}
 B^{-1}:D_s^2\log a
 &-(n+1)(D_s\log a)^{\mathsf T}B^{-1}D_s\log a\\
 {}+R^{-1}:\nabla_S^2\log a
 &+(\nabla_S\log a)^{\mathsf T}R^{-1}\nabla_S\log a
 +\tr R^{-1}=0.
 \end{split}
\end{equation}
Introduce
\[
 \begin{gathered}
 p=D_s\log h,\qquad d=D_sf,
 \qquad p_e=\nabla_S\log h,\qquad d_e=\nabla_Sf,\\
 T_sg=hB^{-1}:D_s^2g,
 \qquad T_eg=hR^{-1}:\nabla_S^2g.
 \end{gathered}
\]
Direct differentiation of $h$ gives
\begin{equation}\label{eq:tube-log-h}
 T_s\log h=-k-|p|^2,
 \qquad
 T_e\log h=n-k-h\tr R^{-1}-|p_e|^2.
\end{equation}
Substitute $\log a=\log h+f$ in
Equation~\eqref{eq:tube-log-euler} and multiply by $h$.  The angular
trace and $|p_e|^2$ terms cancel, and
Equation~\eqref{eq:tube-log-h} gives
\begin{equation}\label{eq:tube-f-equation}
 n-2k+T_sf-|p|^2-(n+1)|p+d|^2
 +T_ef+2\langle p_e,d_e\rangle+|d_e|^2=0.
\end{equation}

\subsection{Weighted integration and coercivity}

Equation~\eqref{eq:tube-f-equation} contains derivatives with different
drifts in the base and angular variables. We integrate with the
invariant measure $\mu$, retaining both drifts until the cancellation
in Equation~\eqref{eq:tube-cancellation}. The dimension restriction
enters only in the quadratic form that remains.

The matrix densities for $T_s$ and $T_e$ relative to
$\dd s\,\dd\omega$ are
\begin{equation}\label{eq:tube-matrix-densities}
 \mu_0hB^{-1}=(\cof B)h^{-n}e^{-cf},
 \qquad
 \mu_0hR^{-1}=(\cof R)h^2e^{cf}.
\end{equation}
For example, the first follows from
$\sqrt{\det R/\det B}=a^{-c}$ and
$h^{1-n/2}a^{-c}=h^{-n}e^{-cf}$; the second follows in the same way
with the signs of the exponents reversed.

Let $v\in C_c^\infty(\mathcal B)$.  The logarithmic derivatives of the
scalar factors in Equation~\eqref{eq:tube-matrix-densities} are
$-np-cd$ and $2p_e+cd_e$, respectively.  One integration by parts and
Equation~\eqref{eq:tube-cofactor-divergence} therefore give, for every
smooth $g$ on $\mathcal T$,
\begin{equation}\label{eq:tube-integration-parts}
 \begin{split}
 \int v^2T_sg\,\dd\mu
 &=\int\left[
   v^2\langle np+cd,D_sg\rangle
   -2v\langle D_sv,D_sg\rangle\right]\dd\mu,\\
 \int v^2T_eg\,\dd\mu
 &=-\int v^2\langle2p_e+cd_e,\nabla_Sg\rangle\,\dd\mu.
 \end{split}
\end{equation}
The angular identity has no cutoff term because $v$ depends only on
$s$.  Taking $g=\log h$ in the first identity and using
Equation~\eqref{eq:tube-log-h} yields
\begin{equation}\label{eq:tube-p-identity}
 \int v^2\left[k+(n+1)|p|^2+c\langle d,p\rangle\right]\dd\mu
 =2\int v\langle D_sv,p\rangle\,\dd\mu.
\end{equation}
Next integrate Equation~\eqref{eq:tube-f-equation} against $v^2$ and
apply Equation~\eqref{eq:tube-integration-parts} with $g=f$.  Expanding
the squares gives
\begin{equation}\label{eq:tube-d-identity}
 \begin{split}
 \int v^2\biggl[\frac n2\bigl(|d|^2+|d_e|^2\bigr)
  +(n+2)\bigl(|p|^2+\langle p,d\rangle\bigr)
  -(n-2k)\biggr]\dd\mu\\
 =-2\int v\langle D_sv,d\rangle\,\dd\mu.
 \end{split}
\end{equation}
In particular, subtracting twice
Equation~\eqref{eq:tube-p-identity} gives the exact cancellation
\begin{equation}\label{eq:tube-cancellation}
 \begin{split}
 \frac n2\int v^2
  \bigl(|d|^2+|d_e|^2-2|p|^2-2\bigr)\,\dd\mu
 ={}&-2\int v\langle D_sv,d\rangle\,\dd\mu\\
    &-4\int v\langle D_sv,p\rangle\,\dd\mu.
 \end{split}
\end{equation}

We combine these identities to control both the $\mu$-mass and the
derivative terms. For a parameter $\lambda>0$, add
$2\lambda/n$ times Equation~\eqref{eq:tube-cancellation} to
Equation~\eqref{eq:tube-p-identity}.  The resulting left integrand is
\[
 F_\lambda=k-2\lambda
 +(n+1-2\lambda)|p|^2+c\langle p,d\rangle
 +\lambda\bigl(|d|^2+|d_e|^2\bigr),
\]
and the identity is
\begin{equation}\label{eq:tube-coercive-identity}
 \int v^2F_\lambda\,\dd\mu
 =\int v\left\langle D_sv,
       \left(2-\frac{8\lambda}{n}\right)p
                -\frac{4\lambda}{n}d\right\rangle\dd\mu.
\end{equation}
In an orthonormal frame for the $s$ covectors, the coefficient matrix
for each paired component of $p,d$ is
\[
 \begin{pmatrix}
 n+1-2\lambda&(n+2)/4\\
 (n+2)/4&\lambda
 \end{pmatrix}.
\]
At $\lambda=1$, its determinant is
\[
 n-1-\frac{(n+2)^2}{16}=\frac{(n-2)(10-n)}{16}.
\]
This is positive for $3\le n\le9$, but the scalar term $k-2$
vanishes when $k=2$. Choosing $\lambda$ slightly below one makes the
scalar term positive while preserving the positive quadratic form.
For the rest of the argument impose $3\le n\le9$ and fix
$\lambda=15/16$, which works throughout this range. Then
$k-2\lambda\ge1/8$, and the determinant satisfies
\[
 \lambda(n+1-2\lambda)-\frac{(n+2)^2}{16}
 =\frac{-8n^2+88n-137}{128}\ge\frac7{128}
 \qquad(3\le n\le9),
\]
where the last inequality follows by concavity of the quadratic and
its endpoint values.  Both diagonal entries are positive.  Thus there
is a constant $\kappa_n>0$ such that
\[
 F_\lambda\ge
 \kappa_n\bigl(1+|p|^2+|d|^2+|d_e|^2\bigr).
\]
Young's inequality absorbs the right side of
Equation~\eqref{eq:tube-coercive-identity}, giving
\begin{equation}\label{eq:tube-energy}
 \int v^2\bigl(1+|p|^2+|d|^2+|d_e|^2\bigr)\,\dd\mu
 \le C_n\int|D_sv|^2\,\dd\mu.
\end{equation}
All constants here depend only on $n$; the estimates are uniform for
$2\le k\le n$.

\subsection{The logarithmic inequality}

For $1\le i\le k$, put $\alpha_i=\log s_i$ and define the measure
\begin{equation}\label{eq:log-measure}
 \dd\sigma=\left(1+\sum_{i=1}^k|D_s\log s_i|^2\right)\dd\mu,
 \qquad |D_s\log s_i|^2=\frac{h(B^{-1})_{ii}}{s_i^2}.
\end{equation}
This measure is invariant under invertible linear changes in $y$ and
positive diagonal changes in $s$: the metric $G$ and its volume are
invariant, while each $\alpha_i$ changes only by an additive constant.

\begin{proposition}\label{prop:log-inequality}
Let $3\le n\le9$ and $2\le k\le n$.  On any tube satisfying the
hypotheses of this section, there is a constant $C_n$, depending only
on $n$, such that every
$\widetilde v\in C_c^\infty(\log\mathcal B)$ satisfies
\begin{equation}\label{eq:log-inequality}
 \int_{\mathcal T}\widetilde v(\log s)^2\,\dd\sigma
 \le C_n\int_{\mathcal T}
       |\nabla\widetilde v(\log s)|_{\R^k}^2\,\dd\sigma.
\end{equation}
Here $\log s=(\log s_1,\ldots,\log s_k)$ and the gradient on the right
is Euclidean.  No completeness of the metric $G$ is required.
\end{proposition}

\begin{proof}
Since $T_s\alpha_i=-|D_s\alpha_i|^2$,
Equation~\eqref{eq:tube-integration-parts} gives
\[
 \begin{split}
 \int v^2|D_s\alpha_i|^2\,\dd\mu
 ={}&-\int v^2\langle np+cd,D_s\alpha_i\rangle\,\dd\mu\\
    &+2\int v\langle D_sv,D_s\alpha_i\rangle\,\dd\mu.
 \end{split}
\]
Young's inequality, followed by Equation~\eqref{eq:tube-energy},
therefore gives
\[
 \int v^2|D_s\alpha_i|^2\,\dd\mu
 \le C_n\int|D_sv|^2\,\dd\mu.
\]
Summing these inequalities and the term controlling $\int v^2\dd\mu$
in Equation~\eqref{eq:tube-energy}, we obtain
\[
 \int v^2\,\dd\sigma\le C_n\int|D_sv|^2\,\dd\mu.
\]
Finally take $v(s)=\widetilde v(\log s)$.  The chain rule and
Cauchy--Schwarz in the positive inner product $hB^{-1}$ imply
\[
 \begin{split}
 |D_sv|^2
 &=\left|\sum_{i=1}^k
       (\partial_i\widetilde v)(\log s)D_s\alpha_i\right|^2\\
 &\le |\nabla\widetilde v(\log s)|_{\R^k}^2
       \sum_{i=1}^k|D_s\alpha_i|^2.
 \end{split}
\]
Equation~\eqref{eq:log-inequality} follows from
Equation~\eqref{eq:log-measure}.  Every integration used a compactly
supported cutoff on the smooth tube, which also proves the last
assertion.
\end{proof}
